\chapter{实验}{Experiments}

\section{实验设置}{Experimental Setup}

实验章节应说明数据集、评价指标、基线方法、实现细节和硬件环境。对于机器学习论文，还应报告随机种子、训练轮数、优化器、学习率、批大小以及模型选择标准。所有设置都应服务于一个明确目标：检验本文方法是否真正解决了绪论中提出的问题。

\subsection{数据集}{Datasets}

数据集描述应包括数据来源、样本数量、训练/验证/测试划分、标签含义和预处理方式。如果使用公开数据集，应引用其来源；如果自行构建数据集，应说明采集、清洗、标注和质量控制流程。

\subsection{评价指标}{Metrics}

评价指标应与研究问题一致。分类任务可以报告准确率、宏平均 F1 或类别平均准确率；检索任务可以报告 Recall@K；生成任务则可能需要自动指标和人工评价结合。若指标存在局限，应在实验分析中说明。

\subsection{基线方法}{Baselines}

基线方法需要覆盖简单基线、经典方法和与本文最接近的方法。若计算资源有限，至少应保证最关键基线在相同数据和相同评价脚本下运行。引用他人论文结果时，应明确其设置是否与本文一致。

\section{主要结果}{Main Results}

表\ref{tab:main-results}展示了实验结果表格的写法。正式论文中应保证每个数字都可以追溯到具体实验设置，并在正文中解释最重要的比较关系。

\begin{table}[htbp]
    \centering
    \caption{主要实验结果示例}
    \label{tab:main-results}
    \begin{tabular}{lcc}
        \toprule
        方法 & 指标A & 指标B \\
        \midrule
        Baseline & 80.0 & 75.0 \\
        Ours & 83.5 & 78.2 \\
        \bottomrule
    \end{tabular}
\end{table}


分析主结果时，可以按照“观察、解释、结论”的顺序组织。观察对应表格中的事实，解释对应可能原因，结论则应谨慎限定在实验支持的范围内。例如，若方法只在两个数据集上提升，就不应声称其在所有场景中都更优。

\section{消融实验}{Ablation Study}

消融实验用于验证方法中关键组件是否真正贡献了性能提升。表\ref{tab:ablation-example}给出了一个占位示例。消融实验的设计应与方法章节中的模块一一对应，否则很难证明改进来自本文提出的核心设计。

\begin{table}[htbp]
    \centering
    \caption{消融实验示例}
    \label{tab:ablation-example}
    \begin{tabular}{lcc}
        \toprule
        设置 & 组件A & 指标 \\
        \midrule
        完整模型 & \ding{51} & 83.5 \\
        去除组件A & \ding{55} & 81.0 \\
        \bottomrule
    \end{tabular}
\end{table}


\section{参数敏感性分析}{Sensitivity Analysis}

当方法包含重要超参数时，可以增加参数敏感性分析。该分析不一定需要穷举所有组合，但应覆盖合理范围，并说明默认参数如何选择。若性能对某个超参数非常敏感，应在局限性中说明这一风险。

\section{误差分析}{Error Analysis}

除总体指标外，还应分析模型在哪些样本、类别或场景下失败。误差分析可以帮助判断方法局限，并为后续改进提供证据。常见做法包括：

\begin{itemize}
    \item 按类别统计错误率，找出困难类别或长尾类别。
    \item 可视化典型成功与失败案例，分析模型依赖了哪些线索。
    \item 比较不同模型在同一批错误样本上的表现，判断错误是否来自数据本身。
\end{itemize}

\section{复现性说明}{Reproducibility}

复现性说明可以放在实验章节末尾。建议列出代码运行入口、依赖环境、随机种子、主要配置文件和模型权重保存方式。对于毕业论文而言，即使不公开完整代码，也应保证自己能够根据论文描述重新运行主要实验。
